\section{Problem Definition and Overview} 
\label{sec:def}
\revised{We begin by formally defining the concept of provenance (Section~\ref{sec:prov-def}) and outlining our cost model (Section~\ref{subsec:cost-model}), followed by a description of our problem statement (Section~\ref{subsec:problem-def}). We further define the applicable tasks (Section~\ref{subsec:strong-weak}) and questions (Section~\ref{subsec:database-prov-relationship}) for which our provenance is useful, and discuss the impact of LLM behavior (e.g., non-determinism) on provenance (Section~\ref{subsec:llm-uncertainty}). 
We finally present an overview of \sys (Section~\ref{subsec:overview}).  }

\vspace{-1mm}
\subsection{Provenance Definition} \label{sec:prov-def}
Consider a data processing task generally represented as a natural language {\em question} $Q$. 
%$Q$ could express, for example, a task that involves 
%data extraction, cleaning, transformation, or summarization.  
Let $L$ be the Large Language Model (LLM), and $T$ be the input {\em text}, 
i.e., source data to process.  
$T$ consists of a sequence of sentences, 
i.e., $T = \langle s_1, s_2, \ldots, s_n \rangle$, 
where $i$ is the {\em index} of sentence $s_i$. 
Let $A$ be the {\em response} (or answer) 
to question $Q$ in text $T$ using LLM $L$, 
denoted as $L(T, Q) = A$. 
Let $P \sqsubseteq T$ be a subsequence of $T$, i.e., 
$P = \langle s_{i_1},s_{i_2},...,s_{i_m} \rangle$, where $1 \leq i_1 < ... < i_m \leq n$. 
%\rone{When a table is part of $T$, it can be converted into text by encoding each tuple as a list of key-value pairs, forming a sentence in $T$, where the key is a column and the value is the corresponding entry in the tuple. }  



\begin{definition}
    [Verifiable Provenance.]~\label{def:provenance} Given text $T$, question $Q$, and LLM $L$, where $L(T, Q) = A$, we define the subsequence $P\sqsubseteq T$ to be a {\em verifiable provenance} if $L(P, Q) = A'$, where $I(A, A') = \text{True}$. 
\end{definition}
\vspace{-3mm}


\noindent The indicator function $I(A, A')$ returns {\em True} if $A$ and $A'$ are lexically identical or semantically equivalent. 
For settings where $A$ 
is drawn from a closed domain (e.g., true/false or zip code), 
users can opt for $I(A, A')$ as $A$=$A'$. Otherwise, 
we leverage an LLM as a judge~\cite{gu2024survey,zheng2023judging} to evaluate
semantic equivalence \revised{(e.g., ``2012'' and ``year 2012''  judged equivalent)}.  
This approach effectively handles equivalent but non-identical LLM-generated responses. \rmix{To show this, we conducted an evaluation on 2000 text pairs from the GLUE benchmark~\cite{wang2018glue}, where each pair contains two text snippets with a label indicating  equivalence. We evaluated LLM-as-a-judge using five LLMs with different capabilities and four prompts, including a handwritten prompt and an LLM-rewritten variant, each with and without examples. \papertext{Details of prompts and experimental setup can be found in ~\cite{blip}}. \techreport{Details of prompts and experimental setup can be found below.} 
As shown in Table~\ref{tab:llm-as-a-judge}, LLM-as-a-judge is effective with around 94\% accuracy on average across models.} 
%We further discuss LLM behavior in Section~\ref{subsec:llm-uncertainty}. 
% \revised{This is effective 
% at tolerating the non-determinism in LLM-generated responses, achieving over 99\% deterministic outputs, as shown in Section~\ref{subsec:llm-uncertainty}, where we further discuss how to handle the rare cases of non-deterministic behaviors in LLMs.  }
\techreport{The LLM used to evaluate $I(A, A')$ does not 
need to be the one used to answer $Q$ on $T$. }


%----


\techreport{\rmix{\topic{\code{response\_equal\_human}}: \textcolor{blue}{[sentence\_1]} and \textcolor{blue}{[sentence\_2]} are placeholders for two test sentences. 
\mypython{
\small 
    llm\_equal\_human: Are the following two sentences semantically equivalent? Please respond with True if they are, and False if they are not. Sentence 1: \textcolor{blue}{[sentence\_1]}. Sentence 2: \textcolor{blue}{[sentence\_2]}. }}}

\techreport{\rmix{\topic{\code{response\_equal\_human\_example}}: \textcolor{blue}{[sentence\_1]} and \textcolor{blue}{[sentence\_2]} are placeholders for two test sentences. 
\mypython{
\small 
    llm\_equal\_human\_example: Are the following two sentences semantically equivalent? Please respond with True if they are, and False if they are not. Examples:
    Example 1:
    Sentence 1: The cat is sleeping on the sofa.
    Sentence 2: A cat is lying on the couch asleep.
    Answer: True
    Example 2:
    Sentence 1: The company reported a profit of 2 million dollars.
    Sentence 2: The company reported a loss of 2 million dollars.
    Answer: False 
Sentence 1: \textcolor{blue}{[sentence\_1]}. Sentence 2: \textcolor{blue}{[sentence\_2]}. }}}

\techreport{\rmix{\topic{\code{response\_equal\_llm}}: \textcolor{blue}{[sentence\_1]} and \textcolor{blue}{[sentence\_2]} are placeholders for two test sentences. 
\mypython{
\small 
    llm\_equal\_2: Determine whether the following two sentences are semantically equivalent. Return True if they are, and False otherwise. Are the following two sentences semantically equivalent? 
Two sentences are semantically equivalent if they express the same meaning, even if the wording is different. 
Minor paraphrasing, reordering, or synonym replacement does NOT change equivalence. 
Differences in factual content, numerical values, negation, or core claims mean they are NOT equivalent. 
Sentence 1: \textcolor{blue}{[sentence\_1]}. Sentence 2: \textcolor{blue}{[sentence\_2]}. } } }

\techreport{\rmix{\topic{\code{response\_equal\_llm\_example}}: \textcolor{blue}{[sentence\_1]} and \textcolor{blue}{[sentence\_2]} are placeholders for two test sentences. 
\mypython{
\small 
    llm\_equal\_2: Determine whether the following two sentences are semantically equivalent. Return True if they are, and False otherwise. Are the following two sentences semantically equivalent? 
Two sentences are semantically equivalent if they express the same meaning, even if the wording is different. 
Minor paraphrasing, reordering, or synonym replacement does NOT change equivalence. 
Differences in factual content, numerical values, negation, or core claims mean they are NOT equivalent. Examples: 

Example 1:
Sentence 1: She did not attend the meeting.
Sentence 2: She skipped the meeting.
Answer: True

Example 2:
Sentence 1: The experiment was conducted in 2020.
Sentence 2: The experiment was conducted in 2021.
Answer: False 
Sentence 1: \textcolor{blue}{[sentence\_1]}. Sentence 2: \textcolor{blue}{[sentence\_2]}. } } }

%----


Intuitively, any subsequence of $T$ that returns (lexically or semantically) equivalent answer to $A$ using the same LLM is considered a verifiable provenance. 
When obvious from  context,  we refer to a
verifiable provenance simply as {\em provenance}. 
In practice, users prefer the provenance to be small (i.e., with fewer sentences) for easier interpretation. Thus, we define the notion of a {\em minimal provenance}. 

\begin{definition}
    [Minimal Provenance.]~\label{def:minimal-provenance} A provenance $P$ is {\em minimal}, if $\forall s\in P, I(L(P\setminus s, Q), A) = False$. 
\end{definition}

\noindent That is, a provenance $P$ is minimal if removing any sentence from $P$ results in an answer that is not equivalent to $A$. 

\begin{definition}
    [Strictly-Minimal Provenance.]~\label{def:strict-minimal-provenance} A provenance $P$ is {\em strictly minimal}, if $\nexists P' \sqsubset P$,  $I(L(P', Q), A) = True$.
\end{definition}

\noindent A provenance $P$ is strictly minimal if there does not exist any subsequence of $P$ that is a provenance. Any strictly minimal provenance is also minimal, by definition. 
In practice, a minimal provenance is often sufficient, as a small-sized, rather than strictly minimal, provenance is typically adequate
for interpreting answers. 
As we will show shortly, under a commonly 
observed assumption for LLMs, 
a minimal provenance is equivalent to a strictly minimal one. 


\begin{table}[bt]
\small 
%\papertext{\color{magenta}}
\centering
\scalebox{0.7}{
\begin{tabular}{l|c|c|c|c|c}
\hline
 & \textbf{gpt-4o} & \textbf{gpt-4o-mini} & \textbf{gpt-5} & \textbf{gpt-5-mini} & \textbf{gemini-2-flash} \\ \hline
{\em response\_equal\_human} & 96.3\% & 91.2\% & 94.2\% & 90.4\% & 94.7\% \\ \hline 
{\em response\_equal\_human\_example} & 96.6\% & 90.8\% & 95.1\% & 90.9\% & 94.3\% \\ \hline 
{\em response\_equal\_llm} & 96.8\% & 90.8\% & 94.7\% & 91.1\% & 95.2\% \\ \hline 
{\em response\_equal\_llm\_example} & 96.4\% & 91.5\% & 93.4\% & 91.6\% & 95.7\% \\ \hline 
\end{tabular}}
\caption{\rmix{\small Effectiveness of LLM-as-a-judge.}} 
\label{tab:llm-as-a-judge}
\vspace{-3em}
\end{table}








\vspace{-1mm}
\subsection{Cost Model and KV-Cache} 
\label{subsec:cost-model}
Let $PT = \langle Q, T \rangle$ denote the prompt formed by concatenating the question $Q$ with the text $T$. We use $L(T, Q)$ and $L(PT)$ interchangeably throughout the paper when it does not lead to confusion. Let $c_{PT}$ represent the monetary cost of a single LLM invocation of $L$ with input $PT$, returning a response $L(PT)$. We have:
\begin{equation}
    c_{PT} = |PT|\cdot c_{in} + |L(PT)|\cdot c_{out}
\end{equation}
\noindent Here, $|\cdot|$ denotes the number of tokens, and $c_{in}$, $c_{out}$ refer to the unit cost of model $L$ for input and output tokens, respectively. For example, when $L$ is \code{gpt-4o-mini}, $c_{in}$ is \$$\frac{0.15}{10^6}$ and $c_{out}$ is \$$\frac{0.6}{10^6}$. 


If two prompts $PT_i$ and $PT_j$ share a prefix, denoted $Pre_{i,j}$, then $Pre_{i,j}$  is considered cached tokens, employing the {\em KV-cache}~\cite{cai2024pyramidkv,liu2024minicache,adnan2024keyformer}. For example, with \code{gpt-4o-mini} and \code{gemini-2-flash}, cached tokens are 2$\times$ and 4$\times$ cheaper than uncached tokens, respectively.
In particular, if the prompt $PT_i$ has already been computed using $L$, then the cost of  inference on $PT_j$, denoted $c_{PT_j \mid PT_i}$, is
\begin{equation}
    c_{PT_j | PT_i} = \frac{|Pre_{i,j}|\cdot c_{in}}{f_L} + |PT_j \setminus Pre_{i,j}|\cdot c_{in} + |L(PT_j)|\cdot c_{out}
\end{equation}
\noindent Here, $f_L$ is the cost reduction factor for cached tokens in model $L$. For example, $f_{\text{gpt-4o-mini}} = 2$, while $f_{\text{gemini-2-flash}} = 4$. 



\begin{table}[bt]
\scriptsize
\centering
\scalebox{0.9}{
\begin{tabular}{c|c||c|c|c|c}
\hline
\textbf{LLMs} & \textbf{Task} & \textbf{Qasper} & \textbf{NL\_DEV} & \textbf{HotpotQA} & \textbf{Average} \\ \hline
\multirow{2}{*}{\textbf{gpt-4o-mini}} 
 & SM & 97\% & 87\% & 88\%  & 90.7\% \\ 
 & WM & 99\% & 96\% & 98\%  & 97.7\% \\ \hline
\multirow{2}{*}{\textbf{gpt-4o}} 
 & SM & 98\% & 93\% & 91\%  & 94.0\% \\ 
 & WM & 100\% & 98\% & 98\%  & 98.7\% \\ \hline
\multirow{2}{*}{\textbf{gemini-2-flash}} 
 & SM & 95\% & 88\% & 81\%  & 88.0\% \\ 
 & WM & 97\% & 96\% & 97\%  & 96.7\% \\ \hline
\end{tabular}}
\caption{\small Percentage of Strongly-Monotonic (SM) and Weakly-Monotonic (WM) Tasks, with per‑row averages.}
\label{tab:tasks}
\vspace{-4em}
\end{table}

\subsection{Problem Definitions}
\label{subsec:problem-def} 

We are now ready to state our problem more formally:
\begin{problem}
    [Provenance Inference.]~\label{def:1-provenance} 
    Given text $T$, question $Q$, LLM $L$ and answer $A=L(T,Q)$, identify a minimal provenance. 
\end{problem}

\noindent The minimal provenance  
in Definition~\ref{def:minimal-provenance} for $\langle Q, A \rangle$ 
within text $T$ may not be unique. 
In our example in Figure~\ref{fig:example}, 
the same officers may be listed multiple times, in the internal affairs report, 
\trs{court hearings report,} as well as incident summary. 
Let $MP = \{mp_1, \dots, mp_n\}$ be the set of distinct minimal provenances within $T$.  
Although in most cases, identifying a single minimal provenance $mp_i \in MP$
is sufficient to provide evidence for human verification,
there are scenarios where multiple minimal provenances
may be helpful, as we will outline in Section~\ref{sec:k-provenance}.
We therefore define the {\em $k$-Provenance Inference} problem below. 
\begin{problem}
    [k-Provenance Inference.]~\label{def:k-provenance} Given text $T$, question $Q$, LLM $L$, and answer $A = L(T,Q)$, identify $k$ distinct minimal provenances,  $\{mp_{i1},mp_{i2},...,mp_{ik}\}$, where $\forall ij, mp_{ij} \in MP$. 
\end{problem}




\vspace{-2mm}
\subsection{Strong and Weak Monotonicity}\label{subsec:strong-weak}
As described previously, our goal, depending on the problem,
is to find one or $k$ minimal provenances.
There is a danger, however, of getting stuck 
at a ``local minimum'', where the provenance $P$ returned is minimal,
but there is a subsequence $P' \sqsubset P$ that is also a provenance---a scenario 
that is addressed by strict minimality.
Here, we introduce two properties and empirically study 
their adherence for our tasks.
We denote a task by the triple $\langle Q,T,L\rangle$.
\begin{definition}
    [Strongly Monotonic Tasks.]~\label{def:strong-mono} A task $\langle Q,T,L\rangle$ with $L(T, Q) = A$, is  {\em strongly monotonic} if for any provenance $P\sqsubset T$, $\forall P'$ where $\forall P\sqsubset P' \sqsubset T$, $P'$ is also a provenance. 
\end{definition}


% \begin{figure}[tb]
%     \centering
%     %\vspace{-1em}
%     \includegraphics[width=0.5\linewidth]{figures/weak-mono.pdf}
%     \vspace{-1em}
%     \caption{\small Weakly Monotonic Tasks. }  
%     \vspace{-1.8em}
%     \label{fig:weak-mono}
% \end{figure}





\noindent A task is {\em strongly monotonic} 
if any superset of any provenance is also a provenance. 
While this may appear restrictive, 
we empirically show in Table~\ref{tab:tasks} 
that strongly monotonic tasks are still common, around 90\% on average
across real workloads with multiple LLMs. (Experimental setup is introduced later.) 
We next introduce the
notion of {\em weakly monotonic} tasks. 

\begin{definition}[Weakly Monotonic Tasks.]
\label{def:weak-mono}
A task $\langle Q,T,L\rangle$ with $L(T,Q)=A$ is {\em weakly monotonic} 
if for any provenance $P \sqsubset T$, where $|P| = i$ and $|T| = n$,
there is a sequence of provenances $P_{i+1}, P_{i+2}, ..., P_{n-1}$
such that $|P_{i+k}| = i+k$, and $P \sqsubset P_{i+1}\sqsubset  P_{i+2}\sqsubset  ...\sqsubset  P_{n-1}\sqsubset T$.
\end{definition}

\begin{figure}[tb]
    \centering
    \includegraphics[width=0.5\linewidth]{figures/weak-mono.pdf}
    \vspace{-1em}
    \caption{\small Weak Monotonicity. }  
    \vspace{-1.8em}
    \label{fig:weak-mono}
\end{figure}
\noindent Instead of requiring that every
superset $P'$ of $P$ is a provenance
as in strong monotonicity, as shown in Figure~\ref{fig:weak-mono}, 
a weakly monotonic task simply requires a {\em chain of provenances}, 
covering every distinct size from $|P|$ to $|T|$, where each larger 
provenance is a superset of smaller ones. 
Weak monotonicity is useful  
because it implies the existence of a {\em path} 
to a minimal 
provenance (e.g., $\langle s_4, s_5\rangle$) 
from  $T$ (e.g., $\langle s_1, \dots, s_9\rangle$). 
We find that {\em weakly monotonic tasks are common,
with a prevalence of {\bf nearly over 95\%} across
workloads and LLMs}. 

We finally connect our notion of weak monotonicity
with the notion of minimal provenance, showing that, for weakly monotonic tasks,
a minimal provenance is also strictly minimal.
\begin{theorem}~\label{theo:minimal-eq}
    For a weakly monotonic task $\langle Q,T,L\rangle$, 
    all minimal provenances $P\sqsubseteq T$ are also strictly minimal. 
\end{theorem}

\noindent \papertext{The proof  
can be found in \vldb{\cite{blip}.}} \techreport{The proof  
can be found below. } 
\rmix{As we will see later, all our solutions for Problem~\ref{def:1-provenance}
are guaranteed to return a minimal provenance for
{\em any} task. } 
Theorem~\ref{theo:minimal-eq} {\bf \em additionally guarantees such minimal provenance is  
also strictly minimal}, 
as long as the task is weakly monotonic--- the dominant
class of tasks in our workload.

\techreport{\begin{proof}
    A strictly minimal provenance is also a minimal provenance by definition. Next, we prove that for a weakly monotonic task $\langle Q, T, L\rangle$, a minimal provenance $P \sqsubseteq T$ is also a strictly minimal provenance. Assume there exists a $P' \sqsubset P$ such that $P'$ is a provenance. We will prove that such a $P'$ does not exist, thereby establishing the correctness of the theorem by contradiction. Let $|P| = n$ and $|P'|=i$. Since $P$ is a minimal provenance, $\forall s \in P$, $P \setminus s$ is not a provenance. Thus, $|P'|$ must be less than $n - 1$. Since the task $\langle Q, T, L\rangle$ is weakly monotonic, there must exist a set of provenances from $P$ to $P'$, i.e., $P' \sqsubset P_{i+1} \sqsubset \dots \sqsubset P_{n-1} \sqsubset P$, where provenance $P_{i+1}$ has size $|i+1|$.  Since $\forall s \in P$, $P \setminus s$ is not a provenance, $\nexists P_{n-1} \sqsubset P$, which contradicts the weak monotonicity property. This completes the proof by contradiction. 
\end{proof}}

%\rtwo{yiming: need to be short and update the flow.} 
\rtwo{So far we have defined strongly and weakly monotonic tasks. We now briefly describe the experimental setup for the results in Table~\ref{tab:tasks}\papertext{, with details in our technical report~\cite{blip}.}\techreport{.} } 

\noindent\rtwo{\textbf{Monotonicity testing protocol.} To compute the percentage of strongly monotonic tasks, for each document-question pair, we begin with a minimal provenance $mp$. We then randomly sample 100 supersets of $mp$, with sizes uniformly distributed between $|mp| + 1$ and $|T| - 1$. We compute the percentage of supersets that reproduce an answer equivalent to $A$, and average this value across all document-question pairs for each dataset. }

\rtwo{To compute the percentage of weakly monotonic tasks, we similarly start from a $mp$. Let $|mp| = i$ and $|T| = n$. Let $\mathcal{P} = \{ P_{i+1}, P_{i+2}, ..., \\ P_{n-1} \}$ be a set of provenances such that $|P_{i+k}| = i+k$, and $P \sqsubset P_{i+1}\sqsubset  P_{i+2}\sqsubset  ...\sqsubset  P_{n-1}\sqsubset T$.  We perform a greedy exhaustive search to construct $\mathcal{P}$. Starting from $mp$, we enumerate candidates for $P_{i+1}$ and stop once one reproduces an answer equivalent to $A$. We then search for $P_{i+2}$, a superset of $P_{i+1}$, and continue this process. If no superset of $P_{i+1}$ with a size of $|i+2|$ yields a valid provenance, we backtrack to explore alternative candidates for $P_{i+1}$.  This continues until $\mathcal{P}$ is discovered\papertext{.} } \techreport{or if not, we treated this as a non-weakly-monotonic task, and we terminate it once the number of LLM-based tests exceeds $m = 500$ to avoid potentially high cost. }

\begin{table}[bt]
\papertext{\vspace{-5pt}}
%\color{blue} 
\scriptsize
\centering
\scalebox{1}{
\begin{tabular}{c|c|c|c|c|c|c}
\hline
& \multicolumn{2}{c|}{\textbf{Qasper}} & \multicolumn{2}{c|}{\textbf{NL\_DEV}} & \multicolumn{2}{c}{\textbf{HotpotQA}}  \\
\hline
& str\_s & llm\_s &str\_s & llm\_s & str\_s & llm\_s \\ \hline 
\textbf{gpt-4o-mini} & 0.94 & \textbf{0.97} & 0.97  & \textbf{0.99}  & 0.99 & \textbf{0.995}  \\ 
\textbf{gpt-4o} & 0.93 & \textbf{0.99}  & 0.97  & \textbf{0.998}  & 0.99  & \textbf{0.998} \\ 
\textbf{gemini-2-flash} & 0.97  & \textbf{0.99}  & 0.96 & \textbf{0.97}  & 0.99 & \textbf{0.99}\\ \hline 
\end{tabular}}
\caption{\small LLM Deterministic Evaluation.}
\label{tab:llm-determininstic}
\vspace{-5em}
\end{table}



%\revised{To ensure determinism on all possible inputs, we extend our provenance definition in Definition~\ref{def:provenance} to handle the special case of output probabilities being tied. Here, we make a simple adjustment when checking if two LLM outputs are equivalent by considering multiple LLM outputs instead of just one. Let $\mathcal{A}$ be the set of all LLM outputs with probability within an adjustment factor $\epsilon$ (a small constant to account for floating point errors) of the one with the maximum probability on the text $T$ and let $\mathcal{A'}$ be the same on a subset of the text $P$. Set $\mathcal{A}$ (or  $\mathcal{A'}$) has more than one element only when output probabilities are tied. Then, to check if $P$ is a verifiable provenance, we can check if $\mathcal{A}$ and $\mathcal{A'}$ share a common element. 
%This change has negligible impact on our results, as tied probabilities occur <1\% of time. } 


\begin{figure*}[tb]
\vspace{-15pt}
\begin{minipage}{0.25\textwidth}
    \centering
    \includegraphics[width=0.9\linewidth]{figures/provenance_correctness.pdf}
    \vspace{-1.3em}
    \caption{\rmix{\small Provenance \\Correctness Taxonomy}}  
    \vspace{-1.3em}
    \label{fig:error_taxonomy}
\end{minipage} 
\begin{minipage}{0.7\textwidth}
    \centering
    \includegraphics[width=0.9\linewidth]{figures/overview2.pdf}
    \vspace{-1.5em}
    \caption{\small Overview of \sys. $T$, $Q$, $A$, and $L$ denote the input text $T$, natural language question $Q$, answer (as evaluated) $A$, and the LLM model used $L$, respectively. $MP$ stands for minimal provenance.}  
    \vspace{-1.5em}
    \label{fig:overview}
\end{minipage} 
\end{figure*}

\vspace{-2mm}
\subsection{Provenance Correctness}
~\label{subsec:database-prov-relationship} 
\rmix{We just discussed minimality properties of verifiable provenance, but have not yet confronted whether the provenance matches user intent. Recall that a verifiable provenance is a subset of the input that is sufficient for an LLM to produce an answer as it would on the whole input. This definition, however, does not imply that the verifiable provenance is sufficient for a human to arrive at the same answer as the LLM.  We call a provenance that enables a human to arrive at the same answer as the LLM a \textit{true provenance}. 
}

 \rmix{We first discuss two scenarios where a true provenance is undefined and thus our provenance is inapplicable.  First, if the question is ambiguous, then the LLM's answer is not meaningful, making the provenance inapplicable. For example, when asking for ``{\em the revenue of Amazon}'' over its financial report, it is unclear whether the question refers to the monthly, quarterly, or yearly revenue. 
Second, if the LLM's answer (on the entire input) is incorrect, then any provenance is incorrect by definition---no provenance is sufficient for a human to arrive at the same answer as the LLM. }


\rmix{When a true provenance is well-defined, a verifiable provenance is \textit{correct} if it matches the true provenance. We next provide a taxonomy of cases where a verifiable provenance can be incorrect (Figure~\ref{fig:error_taxonomy}). Since the true provenance is unknown, this taxonomy provides a way to understand when to apply our techniques.} 


\rmix{\textit{Unstable provenance.} We say that our verifiable provenance is {\em unstable} when it is not meaningful for a given task, so the returned provenance can be incorrect even if the LLM answer is correct. We now discuss set of questions $Q$ for which verifiable provenance is stable.}    As in classical provenance literature~\cite{green2007provenance,buneman2001and,cheney2009provenance,glavic2021data},  $Q$
can represent the computation corresponding to 
any {\em positive relational algebra} (${\mathcal RA}^+$); this class 
encapsulates SPJRU (Select-Project-Join-Rename-Union) queries\trs{, 
along with recursion}. This broad class of queries
encompasses data extraction (e.g., ``extract all
the dates in this text'') or transformation 
(e.g., ``reformat the intake form as a table''). 
$Q$ could also include multi-hop reasoning (e.g., ``what was the origin of
the money that was given to the officer?), analogous to 
a recursive query that follows exchanges of money to its source.
Prior work on {\em why} and {\em where}~\cite{buneman2001and} and {\em how}~\cite{green2007provenance} provenance
also focus on ${\mathcal RA}^+$. 
In fact, Glavic~\cite{glavic2021data} describes
that why provenance is synonymous with ${\mathcal RA}^+$. 
$Q$ could also include as part of its output computation
corresponding to monotonic aggregates, i.e., where
the aggregates trend in a single direction, such as 
\texttt{COUNT(*)} 
or \texttt{SUM} of non-zero positive numbers; e.g., ``count the number of officers involved''
or ``summarize all medical aspects in the document''. 

As in traditional provenance, as soon as $Q$ includes
arbitrary negation or aggregation, or operations on aggregates, 
the usefulness 
of a minimal provenance decreases.
As an extreme example, if $Q$ represents the computation of 
the average of numbers $ \in  \{1, 2, 3, 4, 5\}$,
one minimal provenance could be $\{3\}$, which is not
very useful. 
Likewise, if $Q$ includes a predicate on an aggregate, e.g., ``find
all officers who were mentioned at most three times'',
a minimal provenance might list 
just one mention for each output officer, 
even if they were listed twice or three times, which may not be useful. 

\rmix{\textit{Spurious provenance.} Finally, we say that a verifiable provenance can be {\em spurious} even for ${\mathcal RA}^+$, when LLMs make incorrect judgments. For example, it may incorrectly judge that a true provenance fails to reproduce the answer on the entire input. }


\if 
There are multiple  reasons why a verifiable provenance can 
While all our solutions for Problem~\ref{def:1-provenance}
are guaranteed to return a minimal provenance,
independent of $Q$,
we now discuss the class of questions $Q$ for which
this provenance is {\em useful},
and connect this to classical provenance literature.

As in classical provenance literature~\cite{green2007provenance,buneman2001and,cheney2009provenance} (see \cite{glavic2021data} for a recent survey), \rone{$Q$
can represent the computation corresponding to 
any {\em positive datalog query}; this class 
encapsulates SPJRU (Select-Project-Join-Rename-Union) queries, 
along with recursion. This broad class of queries
encompasses data extraction (e.g., ``extract all
the dates in this text'') or transformation 
(e.g., ``reformat the intake form as a table''). 
$Q$ could also include multi-hop reasoning (e.g., ``what was the origin of
the money that was given to the officer?), analogous to 
a recursive query that follows exchanges of money to its source.} 
Prior work on {\em why} and {\em where}~\cite{buneman2001and} and {\em how}~\cite{green2007provenance} provenance
also focus on this subclass of queries. 
In fact, Glavic~\cite{glavic2021data} describes
that why provenance is synonymous with 
positive relational algebra ${\mathcal RA}^+$,
a class equivalent to SPJRU queries.
 
The question $Q$ could also include as part of its output computation
corresponding to monotonic aggregates, i.e., where
the aggregates trend in a single direction, such as 
\texttt{COUNT(*)} 
or \texttt{SUM} of non-zero positive numbers; e.g., ``count the number of officers involved''
or ``summarize all medical aspects in the document''. 


As in traditional provenance, as soon as $Q$ includes
arbitrary negation or aggregation, or operations on aggregates, 
the usefulness
of a minimal provenance decreases. 
As an extreme example, if $Q$ represents the computation of 
the average of numbers $ \in  \{1, 2, 3, 4, 5\}$,
one minimal provenance could be $\{3\}$, which is not
very useful. 
Likewise, if $Q$ includes a predicate on an aggregate, e.g., ``find
all officers who were mentioned at most three times'',
a minimal provenance might list 
just one mention for each output officer, 
even if they were listed twice or three times, which may not be useful.

Finally, it is worth mentioning that prior work on {\em why} provenance~\cite{buneman2001and,cui2000tracing}
introduces the notion of {\em witnesses}, much like our notion of
provenance, with {\em sufficiency} mirroring our notion of  
{\em verifiability}, while minimality being defined identically.
For ${\mathcal RA}^+$, a property analogous to strong monotonicity holds
for traditional data provenance, where any superset of a provenance is also one;
in our case, we also include the LLM $L$ in our 
definition of the property, since
the behavior can vary by LLM (See Table~\ref{tab:tasks}),
and further relax it to give weak monotonicity, a novel property.
% How provenance is not applicable 
% as the LLM computation is opaque to us. 
\fi

\vspace{-4mm}
\subsection{Handling Non-Determinism in LLM}~\label{subsec:llm-uncertainty}
\revised{So far, we have defined our problems, the applicable tasks based on monotonicity, and the class of questions for which provenance is useful. We now examine LLM behavior and its impact on provenance.} 
When answering questions, instead of producing a single output,  
LLMs produce  
outputs with a probability distribution, which might introduce uncertainty when the same input is processed across multiple runs. While this uncertainty is already handled by our use of LLM-as-a-judge to accomodate for non-identical, but equivalent answers, 
here we go further to see how LLM non-determinism can impact provenance. 

Overall, the degree of uncertainty in LLM output depends on the decoding strategy. 
\rtwo{When greedy decoding is used (by setting the temperature to 0), the model picks the most probable token at each step~\cite{decoding}, and is thus theoretically deterministic. }
%In question answering, since users
%prefer deterministic outputs, greedy
%decoding is used, with temperature as 0,
%and the most probable token picked at each step~\cite{decoding}. 
To validate this, we run each of 500 question–text pairs from our three workloads 50 times. 
\revised{We then cluster the 50 outputs per pair using exact string match (resp., LLM-as-a-judge), and report the maximum cluster size divided by 50 as {\em str\_s} (resp., {\em llm\_s}) in Table~\ref{tab:llm-determininstic}. }
Non-deterministic outputs measured by string match happen about 1\% to 7\% of the time across workloads. 
\rtwo{Such a gap occurs when probabilities are tied}, due to factors like \rtwo{non-deterministic GPU kernels, quantization rounding, or floating-point errors~\cite{floating}}. These cases are 
caught by LLM-as-a-judge, resulting in {\bf deterministic outputs $>$99\%} of the time---so the impact of non-determinism is insignificant. 

\techreport{To ensure determinism on all possible inputs, we  further extend the indicator function $I()$ (in Definition~\ref{def:provenance}) that checks if two LLM outputs are equivalent to handle the special case of output probabilities being tied. Here, instead of considering a single LLM output, let $\mathcal{A}$ and $\mathcal{A'}$ be the set of all LLM outputs on the full text $T$ and a subset of the text $P$. 
\revised{Then, to verify whether $P$ is a verifiable provenance, we check if the outputs with the maximum probability in $\mathcal{A}$ and $\mathcal{A'}$ are equivalent. We introduce a small constant ($\epsilon=0.01$ in our implementation) to account for floating-point errors when identifying outputs with the maximum probability. For example, consider $\mathcal{A} = \{a_1, a_2, a_3\}$ (where $a_i$ is a possible output) with corresponding probabilities  $0.43$, $0.425$, and $0.145$. In this case, both $a_1$ and $a_2$ are outputs with maximum probability under $\epsilon$-relaxation. $\mathcal{A}$ is equivalent to $\mathcal{A'} = \{a'_1, a'_2\}$ (where $a'_1$ has a probability of 0.95) if either $I(a_1, a'_1)$ or $I(a_2, a'_1)$ returns True. } In the following, we adopt this implementation of $I()$, which, however, has a negligible impact on our results, occurring in less than 1\% of cases, compared to the original $I()$ implementation. }



%To ensure determinism on all possible inputs, we  can extend our approach to handle the special case of output probabilities being tied. Here, we make a simple adjustment when checking if two LLM outputs are equivalent by considering multiple LLM outputs instead of just one. Let $\mathcal{A}$ be the set of all LLM outputs with probability within an adjustment factor $\epsilon$ (a small constant to account for floating point errors) of the one with the maximum probability on the full text $T$ and let $\mathcal{A'}$ be the same on a subset of the text $P$. Set $\mathcal{A}$ (or  $\mathcal{A'}$) has more than one element only when output probabilities are tied. Then, to check if $P$ is a verifiable provenance, we can check if $\mathcal{A}$ and $\mathcal{A'}$ share a common element. 
 

 